\section{Objetivos do projeto}
\subsection{Objetivos gerais}
    Investigar e demonstrar como utilizar Modelos de Linguagem de Grande Porte (LLMs) para a geração automática de casos de teste unitários em Java, integrados ao framework JUnit, de modo a satisfazer múltiplos critérios avançados de cobertura de código-fonte (all-paths, all-defs-uses, all-uses)

\subsection{Objetivos específicos}

\begin{itemize}
    \item Construir um pipeline que, dado o código-fonte de um método em Java, gere automaticamente casos de teste JUnit por meio de LLM.
    \item Avaliar a cobertura dos testes gerados usando ferramentas como JaCoCo (ou Cobertura/EclEmma), medindo cobertura de instruções, ramos (branches), definições-uso e usos.
    \item Comparar cobertura e qualidade dos testes gerados via LLM com testes gerados por ferramentas tradicionais (EvoSuite) e com testes manuais, identificando pontos fortes e limitações.
\end{itemize}